\section{System Overview}
\label{sec:sysoverview}
In this section, we presents the system architecture and the lifecycle of RECAS.

\subsection{System Architecture}
%Figure \ref{fig:overview}gives the high-level architecture of RECAS, which consists of two parts: Reconfiguration and Caching. The former is an offline rewrite — rewritten once, and every subsequent reader benefits. The latter is an online cache that performs per-request admission and eviction at runtime, eliminating redundant remote HTTP GETs. Both parts rely solely on evidence observable from the SDK: the byte-level range trace of each request, and the Parquet footer. We briefly introduce the two modules below.
RECAS is a drop-in layer around the object-storage client. It exposes the same \textit{getobject} interface as the underlying SDK, so engines and data loaders need no code changes, and it produces only standard Parquet on the write path. All decisions use evidence the SDK can observe itself --- the byte-level Range GET and the Parquet footer. As shown in Figure~\ref{fig:overview}, RECAS has two modules: \emph{layout reconfiguration} on the write path and \emph{caching} on the read path. These two parts together realize the closed loop motivated earlier: the read path records how data is consumed while the write path rewrites the layout accordingly; subsequent reads then issue fewer remote GETs, which the cache can further eliminate. 

\begin{figure}[htbp]
    \centering
    \includegraphics[width=0.5\textwidth]{figs/overview.png}
    \caption{System Architecture}
    \label{fig:overview}
\end{figure}

%\textbf{Reconfiguration}. The SDK intercepts the headers of outgoing HTTP GET requests to obtain the byte-level access trace, and parses the physical layout of each file from its Parquet footer. The trace and the layout are then correlated geometrically to produce an access profile. The action space for re-laying out a file has five dimensions, and a bounded set of candidates is generated for each dimension; a candidate plan is a point in the Cartesian product of the five dimensions. To ensure that the generated plan is actually writable by the writer, and to compensate for the blind spots of the pricing model, every candidate must first pass an L0 feasibility check, which prunes candidates that violate the hard criteria. The survivors then go through L1 pricing: L1 scales a candidate geometry out of the measured baseline geometry, then prices and ranks every candidate under the objective function. Guided by the plan selected by pricing, a Parquet writer rewrites the original files into the Reconfigured Layout Parquet. From then on, any reader can read this Parquet and benefit from it.
\textbf{Layout Reconfiguration} Outgoing GET headers yield a byte-range access trace; Parquet footers yield the physical geometry of each file. RECAS correlates the two purely geometrically to build an access profile (which columns are read, and which are read together), then searches a five dimensional writer-controlled action space: target file size, row-group size, column order, codec and encoding, and page geometry. A bounded candidate set is generated for each dimension, and each plan is a point in their Catesian product. Every candidate must first pass an \textit{L0} feasibility check, which prunes plans that the writer cannot produce or that violate hard criteria such as a parallelism floor on the number of the scan units, and which also compensates for blind spots of the later pricing model. %Each candidate is first checked for L0 feasibilty (including a parallelism floor on the number of scan units), then priced on a \textit{virtual footer} by replaying the reader's merge rules against the measured baseline --- no data is rewritten to evaluate a point. 
Survivors then go through \emph{L1} pricing: L1 materializes each candidate as a \emph{virtual footer}. scales its geometry relative to the measured baseline, and prices and ranks it by replaying the reader's merge rules --- so no data is rewritten to evaluate a point. The selected plan is materialized through a standard Parquet writer; every subsequent reader benefits without modification. We detail this pipeline in Section IV-A. 

%\textbf{Caching}. In this part, a proxy intercepts \texttt{getObject} and obtains the requested range before each HTTP GET is issued. RECAS normalizes it into a half-open interval \texttt{[start, endExclusive)}, which is used to test whether the content is already cached. On a cache hit, the covering blocks in the cache are pinned and returned as an in-memory stream via zero copy, and no HTTP GET is issued. A cache miss goes through a two-step decision. A size-and-mode filter is applied first: requests that do not pass are served by an exact-range GET, while requests that pass are served from the S3 response stream and their bytes are simultaneously collected into a buffer. Once the buffer is full, a cost-aware  W-TinyLFU algorithm performs request-level atomic cache admission and eviction. Rejected requests are recorded in the range ghost, while admitted requests displace the victims and enter the cache.
\textbf{Caching and Request I/O} A JDK dynamic proxy intercepts synchronous \texttt{getObject} calls, extracts the requested range, and normalizes it to a half-open interval $[start, endExclusive]$, which is used to test whether the content is already cached. On a full-cover hit, RECAS pins the covering blocks and returns an in-memory stream; no HTTP GET is issued. On a miss, a size-and-mode filter first separates the stream: oversized or bypass-mode requests (typically one-shot large scans) are passed through as exact-range GETs and never enter the cache; the remaining requests are served from the S3 response stream while their bytes are teed into request-sized buffer. When the buffer is full, a cost-aware W-TinyLFU policy admits or rejects the request as a whole by frequency $\times$ estimated remote cost. Rejected requests are recorded in a ghost table; admitted ones displace victims and enter the cache. We detail the request path adn admission policy in Section~IV-B.

\subsection{Lifecycle and Coordination}
%Concretely, the system operates
RECAS run in three phases: \textbf{observe}, \textbf{reconfigure}, and \textbf{serve}. 
%In the observe phase, a single baseline run records the offset, length, start and end timestamps, and thread of every request, which feeds candidate pricing in the reconfigure phase. In the reconfigure phase, run offline exactly once, the layout is rewritten through the Reconfiguration process described above, and the output is standard Parquet. In the serve phase, activated at runtime in each process, the transparent cache is turned on and performs per-request admission and eviction at runtime, eliminating redundant remote range GETs. The two parts are weakly dependent by design — re-laying out columns does not change the cache's admission logic, and the cache does not change the footer geometry — so each can be deployed on its own, or the two can be stacked. The 2×2 grid in §x is at once the final result and an empirical test of this weak-dependence assumption. Note that the collector is independent of the cache layer: the trace must be collected in a baseline run with the cache disabled, since hits would otherwise suppress GETs and the trace would no longer reflect the reader's true byte demand.

\begin{itemize}
    \item \textbf{Observe.} Run once as a baseline. It records every request's offset, length, start and end timestamps, and issuing thread, with the cache disabled. This trace feeds candidate pricing in the reconfigure phase. The collector is independent of the cache layer: the trace must be gathered with the cache disabled, because hits would suppress GETs and the trace would no longer reflect the reader's true byte demand.
    \item \textbf{Reconfigure.} Run offline exactly once. This process consumes the observed trace together with the Parquet footers, selects a layout plan, and rewrites the data layout. The output is standard Parquet.
    \item \textbf{Serve.} Activated at the runtime in each process. The transparent cache is turned on and performs per-request admission and eviction, eliminating redundant remote Range GETs. Since the two modules are only weakly coupled, serve can run on the reconfigured layout, or on the canonical layout if reconfiguration is skipped.
\end{itemize}

These three phases coordinate two mechanisms: layout rewriting in \textbf{Reconfigure}, and request-path caching in \textbf{Serve}, with \textbf{Observe} supplying the trace data that pricing the needs. We note that the two mechanisms are wearkly coupled by design: rewriting the layout does not alter the cache's admission logic, and the cache does not alter footer geometry. Either mechanism can thereby be used on its own, or both can be stacked. In the latter case, serve runs on the reconfigured layout; if reconfigure is skipped, it runs on the canonical layout. The $2\times2$ factorial evaluation (i.e. layout% $\in$ \{canonical, reconfigured\} 
$\times$ cache%$\in$ \{off, on\}
, in section~\ref{}), is both the end-to-end result and an empirical test of this weak-coupling assumption.% We detail the design of each mechanism in Section IV.